\section{The obstruction from the median plane}
\label{sec:obstruction}

We now fix the density and angle; no parameter tends to zero in this
section.  We use the following analytic properties established above:
\(n>0\) is continuous, even under reflection in \(z=0\), and axisymmetric,
with \(\int_{\R^3}n=3\).  We retain the regularity
\(n\in H^2(\R^3)\), \(\sqrt n\in H^1(\R^3)\), and
\(\Delta n\in L^1(\R^3)\) from \Cref{lem:density-regularity}.
The angle \(\theta\) is odd, locally \(C^1\), and
\(0<\theta<\pi/2\) on \(z>0\).  It has finite weighted gradient energy and
satisfies
\begin{equation}
 \operatorname{div}(n\nabla\theta)=-d n\sin(2\theta),
 \qquad d>0.
 \label{eq:obstruction-angle}
\end{equation}
The orbitals and the formal potential have the properties
\begin{gather}
 g=\sqrt n\cos\theta>0,\qquad h=\sqrt n\sin\theta,
 \qquad \|g\|_2^2=2,\quad \|h\|_2^2=1,
 \label{eq:obstruction-normalization}\\
 g,h\in H^1(\R^3)\cap H^2_{\mathrm{loc}}(\R^3),\qquad
 V=\frac{\Delta g}{2g}\in L^2_{\mathrm{loc}}(\R^3),
 \qquad \int_{\R^3}n|V|<\infty,
 \label{eq:obstruction-regularity}\\
 \left(-\frac12\Delta+V\right)g=0,\qquad
 \left(-\frac12\Delta+V\right)h=dh
 \quad\text{in distributions}.
 \label{eq:obstruction-orbitals}
\end{gather}
In particular, \(g\) and \(V\) are even and \(h\) is odd under the
reflection.  These are the conclusions of
\Cref{prop:orbital-equations,prop:ensemble-certificate}; the arguments below
use only the displayed analytic properties, rather than a small-coupling
hypothesis.  The density translation limits and the planar lower bound
will be invoked with fixed \(k>0\), \(\beta\in\R\), and
\(A(r)=e^{-kr}r^\beta\).

Suppose provisionally that
\begin{equation}
 V\in L^{3/2}(\R^3)+L^\infty(\R^3).
 \label{eq:admissible-formal-potential}
\end{equation}
Then the quadratic form
\[
 q_V[u]=\frac12\int_{\R^3}|\nabla u|^2+
              \int_{\R^3}V|u|^2,
 \qquad u\in H^1(\R^3),
\]
is closed and bounded below.  Indeed, truncating the \(L^{3/2}\) part
leaves an arbitrarily small \(L^{3/2}\) norm, whose form is controlled by
Sobolev's inequality; the truncated part and the original bounded part
are bounded multipliers.  The same estimate shows that the distributional
identities in \Cref{eq:obstruction-orbitals} extend, by density, to all
\(H^1\) form tests.  Thus \(g\) and \(h\) are eigenfunctions of this
ordinary form operator at energies \(0\) and \(d\), respectively.

\subsection{Limits of translated potentials and orbitals}

\begin{lemma}[Compactness under escaping translations]
\label{lem:translated-orbital-compactness}
Write \(V=V_{\mathrm c}+V_{\mathrm b}\), where
\(V_{\mathrm c}\in L^{3/2}(\R^3)\) and
\(V_{\mathrm b}\in L^\infty(\R^3)\).  If \(|y_j|\to\infty\), then
\[
 V_{\mathrm c}(y_j+\cdot)\longrightarrow0
 \quad\text{in }L^{3/2}_{\mathrm{loc}}(\R^3).
\]
After passing to a subsequence there is a real bounded function \(W\) such
that
\[
 V_{\mathrm b}(y_j+\cdot)\stackrel{*}{\rightharpoonup}W
 \quad\text{locally in }L^\infty,
 \qquad \|W\|_\infty\le\|V_{\mathrm b}\|_\infty.
\]
Suppose, in addition, that \(u_j\in H^1_{\mathrm{loc}}\) solves
\[
 \left(-\frac12\Delta+V(y_j+\cdot)\right)u_j=\varepsilon u_j
\]
for one fixed real \(\varepsilon\), and that \(u_j\) is bounded in
\(L^2\) on every compact set.  A further subsequence converges strongly
in \(L^2_{\mathrm{loc}}\) and weakly in \(H^1_{\mathrm{loc}}\) to a
solution of
\[
 \left(-\frac12\Delta+W\right)u=\varepsilon u.
\]
\end{lemma}

\begin{proof}
Every fixed translated ball eventually lies in an arbitrarily distant
tail of \(V_{\mathrm c}\), proving its norm convergence.  Weak-star
compactness and a diagonal subsequence on expanding balls give \(W\).

To prove the assertion about solutions, let \(\chi\) be a smooth compact
cutoff.  Testing the equation with \(\chi^2u_j\) and taking the real part
gives the usual cutoff identity.  Its possibly negative critical-potential
term is bounded by
\[
 \int |V_{\mathrm c}(y_j+x)|\,|\chi u_j|^2\,dx
 \le C_{\mathrm S}\|V_{\mathrm c}(y_j+\cdot)\|_{L^{3/2}(\supp\chi)}
          \|\nabla(\chi u_j)\|_2^2.
\]
For large \(j\) this is absorbed into the kinetic term.  The bounded
potential, the eigenvalue, and the cutoff derivative are controlled by
the assumed local \(L^2\) bounds.  Hence the sequence is bounded in local
\(H^1\).  Rellich compactness gives the stated convergences.

For a compact smooth test \(\varphi\), the critical-potential term tends
to zero by H\"older's and Sobolev's inequalities.  The bounded-potential
term passes to the limit because \(u_j\varphi\to u\varphi\) in \(L^1\)
and the multipliers are uniformly bounded, while their pairing with the
fixed function \(u\varphi\) converges weak star.  The gradient term passes
by weak convergence.  This proves the limiting equation.
\end{proof}

\begin{proposition}[The unique potential limit at infinity]
\label{prop:potential-limits}
Assume \Cref{eq:admissible-formal-potential} and the density translation
limits \Cref{eq:plane-translate,eq:bulk-translate}.  For every sequence
\(|y_j|\to\infty\), every local weak-star subsequential limit of
\(V_{\mathrm b}(y_j+\cdot)\) is the constant
\begin{equation}
 c_\infty=d+\frac{k^2}{2}.
 \label{eq:potential-limit-constant}
\end{equation}
\end{proposition}

\begin{proof}
\textbf{Centers on the median plane.}
First let \(y_j=p_j\omega_j\), where \(p_j\to\infty\),
\(\omega_j\cdot e_z=0\), \(|\omega_j|=1\), and, after subselection,
\(\omega_j\to\omega_*\).  Put
\[
 c_j=\frac{A(p_j)}{p_j},\qquad
 g_j(x)=\frac{g(y_j+x)}{c_j},\qquad
 h_j(x)=\frac{h(y_j+x)}{c_j}.
\]
By \Cref{eq:plane-translate}, for a constant \(a_*>0\),
\begin{equation}
 g_j^2+h_j^2=\frac{n(y_j+\cdot)}{c_j^2}
 \longrightarrow a_*^2z^2e^{-2k\omega_*\cdot x}
 \quad\text{locally uniformly}.
 \label{eq:plane-normalized-orbitals}
\end{equation}
In particular these orbitals are locally bounded.  Take any subsequential
bounded-potential limit \(W\).  By
\Cref{lem:translated-orbital-compactness}, a further subsequence has limits
\(g_\infty,h_\infty\) satisfying the two orbital equations with potential
\(W\) and energies \(0,d\).  Their strong local \(L^2\) convergence and
\Cref{eq:plane-normalized-orbitals} give
\[
 0\le g_\infty(x)\le a_*|z|e^{-k\omega_*\cdot x},\qquad
 g_\infty^2+h_\infty^2=a_*^2z^2e^{-2k\omega_*\cdot x}
 \quad\text{a.e.}
\]
Since \(W\) is bounded, local elliptic regularity makes \(g_\infty\)
continuous.  It vanishes on the interior plane \(z=0\).  The local
Harnack inequality for a nonnegative solution with bounded zeroth-order
coefficient \cite[Lecture~18, Theorem~1]{SimonPDE} forces
\(g_\infty\equiv0\); no sign condition on \(W\) is needed.
The sign of \(h\) on each half-space therefore determines
\[
 h_\infty(x)=a_*z e^{-k\omega_*\cdot x}.
\]
Because \(\omega_*\cdot e_z=0\), its Laplacian is \(k^2h_\infty\).
The limiting equation gives
\((W-d-k^2/2)h_\infty=0\).  This determines \(W=c_\infty\) almost
everywhere; the plane where \(h_\infty=0\) has zero volume.

\smallskip\noindent\textbf{A positive angle bound above the plane.}
We next show
\begin{equation}
 \inf_{x_1,x_2\in\R}\theta(x_1,x_2,1)>0.
 \label{eq:angle-plane-parallel-lower-bound}
\end{equation}
On each compact subset of the open upper half-space, the normalized
density in \Cref{eq:plane-normalized-orbitals} is uniformly bounded above
and below by positive constants.  The translated angle solves
\[
 \operatorname{div}\!\left(\frac{n(y_j+x)}{c_j^2}
                         \nabla\theta(y_j+x)\right)
 =-d\frac{n(y_j+x)}{c_j^2}\sin\bigl(2\theta(y_j+x)\bigr).
\]
On each such compact set the scalar coefficients have uniform positive
lower and finite upper bounds, and the right-hand sides are uniformly
bounded.  Local energy estimates and the interior H\"older estimate for
uniformly elliptic divergence-form equations
\cite[Lecture~18, Theorem~3]{SimonPDE} therefore give local uniform
compactness of the bounded angles.  These estimates require no bounds
on derivatives of the normalized densities.  The identity
\[
 g_j=\sqrt{n(y_j+\cdot)/c_j^2}\cos\theta(y_j+\cdot),
\]
together with \(g_\infty=0\), implies that every such angle limit
equals \(\pi/2\).
Consequently \(\theta(x_1,x_2,1)\to\pi/2\) as
\(x_1^2+x_2^2\to\infty\).  On the remaining compact part it has a
strictly positive minimum, proving
\Cref{eq:angle-plane-parallel-lower-bound}.

Choose \(m\in(0,\pi/2)\) below this infimum.  In \(z>1\), the angle
is a supersolution of \(-\operatorname{div}(n\nabla\theta)\ge0\).
For \(w=(m-\theta)_+\), whose trace on \(z=1\) is zero, test with
\(\chi_R^2w\), where \(\chi_R\) is a spatial cutoff with
\(|\nabla\chi_R|\le C/R\).  The resulting estimate is
\[
 \int_{z>1} n\chi_R^2|\nabla w|^2
 \le4\int_{z>1}n w^2|\nabla\chi_R|^2\longrightarrow0.
\]
Here \(w\) is bounded and \(n\) is integrable.  Thus \(w\) is constant
on this connected half-space, and its zero trace makes it zero.  We have
proved \(\theta\ge m\) throughout \(z\ge1\).

\smallskip\noindent\textbf{Centers whose heights diverge.}
Set \(\delta=\pi/2-\theta\) on the upper half-space.  On \(z\ge1\),
\[
 0\le\delta\le\frac\pi2-m,\qquad
 \operatorname{div}(n\nabla\delta)=d n\sin(2\delta),\qquad
 d\sin(2\delta)\ge c_0\delta
\]
for some \(c_0>0\).  The last assertion follows by minimizing the
continuous positive function \(d\sin(2s)/s\) on
\([0,\pi/2-m]\), with its value at zero defined as \(2d\).

Now let \((y_j)_3\to+\infty\), and pass to a subsequence with
\(y_j/|y_j|\to\omega_*\).  By \Cref{eq:bulk-translate},
\[
 N_j(x)=\frac{n(y_j+x)}{n(y_j)}
 \longrightarrow N_\infty(x)=e^{-2k\omega_*\cdot x}
 \quad\text{locally uniformly}.
\]
The same local energy and H\"older estimates give, along a subsequence,
local uniform and weak local \(H^1\) convergence of
\(\delta(y_j+\cdot)\) to a bounded nonnegative function \(w\).
Passing the divergence equation to the limit requires only the displayed
uniform convergence of the coefficients, not convergence of their
derivatives.  We obtain
\[
 \operatorname{div}(N_\infty\nabla w)=dN_\infty\sin(2w),\qquad
 (\Delta-2k\omega_*\cdot\nabla)w=d\sin(2w)\ge c_0w.
\]
Local elliptic regularity makes \(w\) smooth.  A short maximum-principle
argument shows that \(w=0\).  Indeed, for \(\varepsilon>0\) subtract
\(\varepsilon\rho(x)\), where
\(\rho(x)=\sqrt{1+|x|^2}\).  The boundedness of \(w\) implies that
\(w-\varepsilon\rho\) attains a global maximum at some \(x_\varepsilon\).
At that point,
\[
 c_0w(x_\varepsilon)
 \le (\Delta-2k\omega_*\cdot\nabla)w(x_\varepsilon)
 \le\varepsilon(3+2k),
\]
since \(|\nabla\rho|\le1\) and \(\Delta\rho\le3\).  Comparing the
maximum with any fixed point \(x_0\) gives
\(w(x_0)\le\varepsilon\rho(x_0)+\varepsilon(3+2k)/c_0\).
Letting \(\varepsilon\downarrow0\) proves the assertion.

It follows that \(\delta(y_j+\cdot)\to0\) locally uniformly along every
such sequence.  Hence
\[
 \frac{h(y_j+x)}{\sqrt{n(y_j)}}
 =\sqrt{N_j(x)}\cos\delta(y_j+x)
 \longrightarrow e^{-k\omega_*\cdot x}
 \quad\text{locally uniformly}.
\]
Apply \Cref{lem:translated-orbital-compactness} to this normalized orbital
along any subsequence on which the bounded potentials converge.  Its
limiting eigenfunction equation again identifies
\(W=d+k^2/2\), since the limiting exponential is everywhere positive.
Reflection treats centers with heights tending to \(-\infty\).
Let \(\mathcal R\) denote reflection in the median plane.
Since \(V\) is even, its arbitrary decomposition satisfies
\[
 V_{\mathrm b}(\mathcal R x)-V_{\mathrm b}(x)
 =V_{\mathrm c}(x)-V_{\mathrm c}(\mathcal R x).
\]
Escaping translates of the right side tend to zero in local
\(L^{3/2}\).  The left side is uniformly bounded, so its local
weak-star limits are also zero.  Reflecting the center and local
coordinates therefore gives the same constant limit for negative heights.

\smallskip\noindent\textbf{All remaining sequences.}
Given any escaping sequence, its heights have a subsequence that tends
to \(+\infty\), tends to \(-\infty\), or converges to a finite number.
The first two cases were just proved.  In the third, project the centers
onto the median plane.  Their horizontal radii tend to infinity, and the
original centers differ from these plane centers by convergent bounded
displacements.  The plane-centered conclusion therefore applies.
To justify the moving displacement for weak-star convergence, translate
each fixed compactly supported \(L^1\) test by the negative displacement;
these tests converge in \(L^1\), while the multipliers stay uniformly
bounded.  Thus their limit remains the same constant.  These cases exhaust
every subsequential bounded-potential limit.
\end{proof}

\subsection{From translated limits to an exterior form bound}

\begin{lemma}[Uniform exterior form convergence]
\label{lem:exterior-form-limit}
Let \(V=V_{\mathrm c}+V_{\mathrm b}\) be as in
\Cref{lem:translated-orbital-compactness}.  Suppose every bounded-potential
limit along escaping translates is a fixed constant \(c_\infty\).
For every \(\varepsilon>0\) there is \(R_\varepsilon<\infty\) such that
\begin{equation}
 q_V[u]\ge
 \left(\frac12-\varepsilon\right)\|\nabla u\|_2^2
 +(c_\infty-\varepsilon)\|u\|_2^2
 \label{eq:exterior-form-coercivity}
\end{equation}
for every \(u\in H^1(\R^3)\) vanishing almost everywhere on
\(\{|x|<R_\varepsilon\}\).
\end{lemma}

\begin{proof}
Fix a ball \(B_L\).  We first claim
\begin{equation}
 \sup_{\substack{u\in H^1_0(B_L)\\\|u\|_{H^1(B_L)}\le1}}
 \left|\int_{B_L}\bigl(V(y+x)-c_\infty\bigr)|u(x)|^2\,dx\right|
 \longrightarrow0\qquad (|y|\to\infty).
 \label{eq:local-form-uniformity}
\end{equation}
Sobolev's inequality and the vanishing local \(L^{3/2}\) norm prove the
claim for the critical part.  If the bounded part failed, choose escaping
centers and tests violating it.  Extract a weak-star subsequence of the
bounded multipliers and a subsequence of the tests converging strongly
in \(L^2(B_L)\).  Their squares converge in \(L^1\), because
\[
 \bigl\||u_j|^2-|u|^2\bigr\|_1
 \le(\|u_j\|_2+\|u\|_2)\|u_j-u\|_2.
\]
Uniform boundedness of the multipliers handles this changing test, and
their weak-star limit is \(c_\infty\).  This contradicts the selected
violation and proves \Cref{eq:local-form-uniformity}.

Choose once and for all a smooth, locally finite quadratic partition
\(\sum_\ell\chi_\ell^2=1\) with uniformly bounded overlap, each
\(\chi_\ell\) supported in a ball of one fixed radius, and
\(\sum_\ell|\nabla\chi_\ell|^2\le K\).  For any prescribed
\(\eta>0\), \Cref{eq:local-form-uniformity} applies with error \(\eta\)
to all these balls whose centers are sufficiently distant.  If \(u\)
is supported sufficiently far out, every nonzero \(\chi_\ell u\) has
such a center.  Summing the local potential estimates gives
\begin{align*}
 \left|\int(V-c_\infty)|u|^2\right|
 &\le\eta\sum_\ell\|\chi_\ell u\|_{H^1}^2\\
 &\le\eta\|\nabla u\|_2^2+\eta(1+K)\|u\|_2^2.
\end{align*}
Here the cross terms cancel because
\(\sum_\ell\chi_\ell\nabla\chi_\ell=0\).  Choosing
\(\eta\le\varepsilon/(1+K)\) proves
\Cref{eq:exterior-form-coercivity}.  This argument uses a uniform local
form estimate, not a pointwise limit of \(V\).
\end{proof}

\subsection{The weighted trace contradiction}

\begin{theorem}[A fixed-density obstruction]
\label{thm:obstruction}
Let \(n,\theta,g,h,V,d\) have the analytic properties stated at the
beginning of this section, with \(d>0\).  Suppose that
\Cref{eq:plane-tail,eq:plane-translate,eq:bulk-translate} hold with fixed
\(k>0\) and \(\beta\), and that the weighted gap
\Cref{eq:weighted-gap} holds.  Then
\[
 V\notin L^{3/2}(\R^3)+L^\infty(\R^3).
\]
Moreover, \(n\) is not the density of a ground-state ensemble of any
three-electron noninteracting Hamiltonian with a single real,
spin-independent potential in this class.  These hypotheses are all
fixed-parameter hypotheses; no perturbative neighborhood is required.
\end{theorem}

\begin{proof}
Assume first that \Cref{eq:admissible-formal-potential} holds.
\Cref{prop:potential-limits,lem:exterior-form-limit} give
\Cref{eq:exterior-form-coercivity} with
\(c_\infty=d+k^2/2\).  Since \(d>0\), choose
\begin{equation}
 k<a<\sqrt{k^2+2d},\qquad
 0<\varepsilon<\min\left\{\frac12,
                       c_\infty-\frac{a^2}{2}\right\}.
 \label{eq:agmon-exponent-choice}
\end{equation}
We use the exponential-weight method \cite{Agmon}, spelling out the
estimate and its \(H^1\) conclusion for the form bound obtained here.

Let \(\eta\) be a smooth cutoff which vanishes on the ball required by
\Cref{eq:exterior-form-coercivity} and equals one outside a slightly larger
ball.  With \(\rho(x)=\sqrt{1+|x|^2}\), put
\[
 f_T(x)=a\min\{\rho(x),T\},\qquad
 u_T=\eta e^{f_T}g.
\]
The function \(f_T\) is bounded and Lipschitz with
\(|\nabla f_T|\le a\).  Thus \(u_T\in H^1\), and the weak eigenfunction
equation for \(g\) may be tested with \(\eta^2e^{2f_T}g\).  Expanding
the gradient gives the exact identity
\begin{equation}
 q_V[u_T]=\frac12\int g^2
                    \bigl|\nabla(\eta e^{f_T})\bigr|^2.
 \label{eq:agmon-cutoff-identity}
\end{equation}
The part with \(\eta^2|\nabla f_T|^2\) is at most
\(a^2\|u_T\|_2^2\).  All terms involving \(\nabla\eta\), including
the cross term, are supported in a fixed compact annulus, where
\(f_T\le a\rho\) independently of \(T\).  Consequently
\[
 q_V[u_T]\le\frac{a^2}{2}\|u_T\|_2^2+C
\]
with \(C\) independent of \(T\).  Combining this with the exterior form
bound yields
\[
 \left(\frac12-\varepsilon\right)\|\nabla u_T\|_2^2+
 \left(c_\infty-\varepsilon-\frac{a^2}{2}\right)\|u_T\|_2^2\le C.
\]
Both coefficients are positive by \Cref{eq:agmon-exponent-choice}.
Monotone convergence of the squared weighted functions gives the
\(L^2\) bound as \(T\to\infty\), and weak \(H^1\) compactness identifies
the limit as \(\eta e^{a\rho}g\).  Adding back the compact interior part
proves
\begin{equation}
 e^{a\rho}g\in H^1(\R^3),\qquad
 e^{a|x|}g\in H^1(\R^3).
 \label{eq:weighted-ground-orbital}
\end{equation}
The second assertion also follows by multiplying the first function by
the bounded Lipschitz factor \(e^{a(|x|-\rho(x))}\).

The trace map from \(H^1(\R^3)\) to \(L^2(\R^2)\) is bounded.
The local continuity of \(g\), \(n\), and \(\theta\), together with the
zero angle trace, identifies the trace on \(z=0\) as
\[
 g(x_1,x_2,0)=\sqrt{n(x_1,x_2,0)}.
\]
Thus \Cref{eq:weighted-ground-orbital} implies
\(\int_{z=0}e^{2ap}n(p,0)<\infty\).  On the other hand,
\Cref{eq:plane-tail} gives a fixed \(c_*>0\) and \(p_*<\infty\) such that
\begin{align*}
 \int_{z=0}e^{2ap}n(p,0)\,dx_1dx_2
 &\ge2\pi c_*^2\int_{p_*}^\infty
            e^{2(a-k)p}p^{2\beta-7}\,dp\\
 &=\infty,
\end{align*}
because \(a>k\).  This contradicts the trace bound and proves that
\Cref{eq:admissible-formal-potential} is impossible.

Finally, the weighted gap and the analytic properties above give the
ensemble kinetic certificate of \Cref{prop:ensemble-certificate}.
If an admissible potential represented \(n\) by a ground-state ensemble,
\Cref{prop:potential-identification} would identify it with \(V\) up to
a constant.  The sum space is invariant under adding constants, so this
would imply \Cref{eq:admissible-formal-potential}, already contradicted.
\end{proof}

\begin{theorem}[A family of finite Coulombic counterexamples]
\label{thm:qualitative-counterexample}
For every sufficiently large fixed \(D\), there is an integer
\(Z_0(D)\) such that, for every integer \(Z\ge Z_0(D)\), the
three-electron Coulomb Hamiltonian with nuclear charges \(Z,Z\) at
\((0,0,\pm D/Z)\) has an absolute ground state whose density is not
representable by a ground-state ensemble of a single real,
spin-independent noninteracting potential in
\(L^{3/2}(\R^3)+L^\infty(\R^3)\).
\end{theorem}

\begin{proof}
First choose \(D\) large enough for
\Cref{prop:two-centre-spectrum,prop:density-properties,prop:exterior-gap}.
For this fixed \(D\), take \(\lambda>0\) sufficiently small that
\Cref{prop:small-coupling,prop:angle-convergence,prop:weighted-gap}
apply as well.  The resulting three-electron state is an absolute bound
ground state, and its density and angle satisfy every hypothesis of
\Cref{thm:obstruction}.  The planar coefficient is nonzero because the
coupling is positive.  All these properties hold throughout some
interval \(0<\lambda<\lambda_0(D)\).  Choose the integer \(Z_0(D)\)
so that \(1/Z_0(D)<\lambda_0(D)\), set \(\lambda=1/Z\), and apply the
scaling equivalence of \Cref{prop:scaling}.
\end{proof}

The parameters in \Cref{thm:qualitative-counterexample} are finite but
are still selected by sufficient-size conditions.  The next section
specifies one charge by an exact spectral and variational formula and
checks its hypotheses directly using the fixed-density form of
\Cref{thm:obstruction}.
